\section{2014 Analysis \& Differential Equations}

\subsection{Individual}

\begin{exercise}
Let $f:\mathbb{R}\to\mathbb{R}$ be a continuous function which satisfies
\[
\sup_{x,y\in\mathbb{R}}|f(x+y)-f(x)-f(y)| < \infty.
\]
If we have $\lim_{n\to\infty,n\in\mathbb{N}}\frac{f(n)}{n} = 2014$, prove that $\sup_{x\in\mathbb{R}}|f(x)-2014x| < \infty$.
\end{exercise}

\begin{solution}
\textbf{Prerequisites:} Cauchy's Functional Equation, Mathematical Induction, Properties of Supremum, and the concept of Additive Functions. \\[1em]
\textbf{Intuition:} The condition $\sup |f(x+y)-f(x)-f(y)| < \infty$ tells us that $f$ is "almost" an additive function (like $f(x)=kx$). If it were exactly additive and continuous, it would be $f(x)=kx$. The problem asks us to show that even with this "error," the function $f$ stays within a bounded distance of a linear function, specifically $2014x$. \\[1em]
\textbf{Steps:}
1. Let $C = \sup_{x,y\in\mathbb{R}}|f(x+y)-f(x)-f(y)|$. This constant $C$ represents the maximum "deviation" from additivity. \\[1em]
2. We first prove by induction that for any $n \in \mathbb{N}$ and $x \in \mathbb{R}$, $|f(nx) - nf(x)| \leq (n-1)C$.
For $n=1$, $|f(x) - f(x)| = 0 \leq 0$.
Assume it holds for $n$: $|f(nx) - nf(x)| \leq (n-1)C$.
For $n+1$:
$|f((n+1)x) - (n+1)f(x)| = |f(nx+x) - f(nx) - f(x) + f(nx) - nf(x)|$
$\leq |f(nx+x) - f(nx) - f(x)| + |f(nx) - nf(x)|$
$\leq C + (n-1)C = nC$.
Thus, the bound holds for all $n$. \\[1em]
3. Define a sequence of functions $A_n(x) = \frac{f(2^n x)}{2^n}$. We can show this sequence converges uniformly to an additive function $A(x)$. Note that $|A_{n+1}(x) - A_n(x)| = \frac{|f(2 \cdot 2^n x) - 2f(2^n x)|}{2^{n+1}} \leq \frac{C}{2^{n+1}}$. This geometric series converges, so $A(x) = \lim_{n\to\infty} A_n(x)$ exists and is continuous because $f$ is continuous and the convergence is uniform. \\[1em]
4. For any continuous additive function $A(x)$, it is well known that $A(x) = kx$ for some $k \in \mathbb{R}$.
Taking the limit in $|f(2^n x) - 2^n f(x)| \leq (2^n-1)C$ divided by $2^n$ as $n\to\infty$, we get $|A(x) - f(x)| \leq C$ for all $x$. \\[1em]
5. Finally, we determine $k$. We are given $\lim_{n\to\infty} \frac{f(n)}{n} = 2014$. Since $f(n) = kn + O(1)$, we have $\frac{kn + O(1)}{n} = k + \frac{O(1)}{n} \to k$. Thus $k = 2014$.
Therefore, $|f(x) - 2014x| \leq C$ for all $x \in \mathbb{R}$, which implies $\sup_{x\in\mathbb{R}}|f(x)-2014x| \leq C < \infty$.
\end{solution}

\begin{exercise}
Let $f_1,\dots,f_n$ be analytic functions on $D=\{|z|<1\}$ and continuous on $\bar{D}$. Prove that $\phi(z)=|f_1(z)|+|f_2(z)|+\cdots+|f_n(z)|$ achieves maximum values at the boundary $\partial D$.
\end{exercise}

\begin{solution}
\textbf{Prerequisites:} Analytic functions, Maximum Modulus Principle, Triangle Inequality, and basic Complex Analysis. \\[1em]
\textbf{Intuition:} The Maximum Modulus Principle states that the absolute value of a non-constant analytic function cannot have a local maximum inside its domain. Although the sum of absolute values $\phi(z)$ is not itself analytic, we can "force" it to look like the modulus of an analytic function at any specific point. This allows us to use the principle to show that the maximum must occur on the boundary. \\[1em]
\textbf{Steps:}
1. Suppose $\phi(z)$ achieves its maximum value at some point $z_0$ in the interior of the disk $D$. We want to show that there is a point on the boundary $\partial D$ where $\phi$ is at least as large. \\[1em]
2. For each function $f_j(z)$, let $\theta_j$ be the phase angle of $f_j(z_0)$. Specifically, we write $f_j(z_0) = |f_j(z_0)|e^{i\theta_j}$. This means $e^{-i\theta_j} f_j(z_0) = |f_j(z_0)|$, which is a real, non-negative number. \\[1em]
3. Define a new function $g(z) = \sum_{j=1}^n e^{-i\theta_j} f_j(z)$. Since each $f_j(z)$ is analytic on $D$ and continuous on $\bar{D}$, their linear combination $g(z)$ is also analytic on $D$ and continuous on $\bar{D}$. \\[1em]
4. Evaluate $g$ at the point $z_0$:
$g(z_0) = \sum_{j=1}^n e^{-i\theta_j} f_j(z_0) = \sum_{j=1}^n |f_j(z_0)| = \phi(z_0)$.
Note that $g(z_0)$ is equal to $|g(z_0)|$ because it is a sum of non-negative real numbers. \\[1em]
5. By the Maximum Modulus Principle, the maximum of $|g(z)|$ on the closed disk $\bar{D}$ must be achieved at some point $z^*$ on the boundary $\partial D$. Thus, $|g(z_0)| \leq |g(z^*)|$. \\[1em]
6. Now we use the Triangle Inequality on the boundary point $z^*$:
$|g(z^*)| = \left| \sum_{j=1}^n e^{-i\theta_j} f_j(z^*) \right| \leq \sum_{j=1}^n |e^{-i\theta_j}| |f_j(z^*)| = \sum_{j=1}^n |f_j(z^*)| = \phi(z^*)$. \\[1em]
7. Putting it all together: $\phi(z_0) = |g(z_0)| \leq |g(z^*)| \leq \phi(z^*)$.
This shows that for any point $z_0$ in the interior, there exists a point $z^*$ on the boundary such that $\phi(z^*) \geq \phi(z_0)$. Therefore, the maximum of $\phi(z)$ must be attained on the boundary $\partial D$.
\end{solution}

\begin{exercise}
Prove that if there is a conformal mapping between the annulus $\{z\mid r_1<|z|<r_2\}$ and the annulus $\{z\mid\rho_1<|z|<\rho_2\}$, then $\frac{r_2}{r_1} = \frac{\rho_2}{\rho_1}$.
\end{exercise}

\begin{solution}
\textbf{Prerequisites:} Conformal mapping, Harmonic functions, Dirichlet integral, and the concept of a Modulus of an annulus. \\[1em]
\textbf{Intuition:} Not all annuli are conformally equivalent. While any two disks can be mapped to each other, an annulus has a "shape" determined by the ratio of its radii. This ratio is a conformal invariant, meaning a conformal map cannot change it. We prove this by looking at a harmonic function that "measures" the distance between the boundaries and showing that its energy (Dirichlet integral) must be preserved. \\[1em]
\textbf{Steps:}
1. Without loss of generality, we can scale the annuli. Let $A_1 = \{z \mid 1 < |z| < R\}$ and $A_2 = \{w \mid 1 < |w| < S\}$, where $R = r_2/r_1$ and $S = \rho_2/\rho_1$. A conformal map $f: A_1 \to A_2$ exists if and only if the original map exists. \\[1em]
2. Consider the harmonic function $u(z) = \frac{\log|z|}{\log R}$ on $A_1$. This function is chosen because it is harmonic ($\Delta u = 0$), equals $0$ on the inner boundary ($|z|=1$), and equals $1$ on the outer boundary ($|z|=R$). \\[1em]
3. Calculate the Dirichlet integral (energy) of $u$:
$D(u) = \iint_{A_1} |\nabla u|^2 dx dy$.
In polar coordinates, $\nabla u = \frac{\partial u}{\partial r} \hat{r} = \frac{1}{r \log R} \hat{r}$.
$D(u) = \int_0^{2\pi} \int_1^R \left( \frac{1}{r \log R} \right)^2 r dr d\theta = \frac{2\pi}{(\log R)^2} \int_1^R \frac{1}{r} dr = \frac{2\pi \log R}{(\log R)^2} = \frac{2\pi}{\log R}$. \\[1em]
4. Since $f$ is a conformal map, the composition $v = u \circ f^{-1}$ is a harmonic function on $A_2$. Because $f$ maps boundaries to boundaries, $v$ will take values $0$ and $1$ on the boundaries of $A_2$ (possibly swapped). The unique such harmonic function on $A_2$ is $v(w) = \frac{\log|w|}{\log S}$ (or $1 - \frac{\log|w|}{\log S}$). \\[1em]
5. A key property of the Dirichlet integral is that it is invariant under conformal mappings. This is because the Jacobian of a conformal map $f(z)$ is $|f'(z)|^2$, which exactly cancels out the transformation of $|\nabla u|^2$.
Thus, $D(u) = D(v)$. \\[1em]
6. From our formula in step 3, $D(v) = \frac{2\pi}{\log S}$.
Setting $D(u) = D(v)$, we get:
$\frac{2\pi}{\log R} = \frac{2\pi}{\log S} \implies \log R = \log S \implies R = S$.
Substituting the ratios back, we have $\frac{r_2}{r_1} = \frac{\rho_2}{\rho_1}$.
\end{solution}

\begin{exercise}
Let $U(\xi)$ be a bounded function on $\mathbb{R}$ with finitely many points of discontinuity. Prove that
\[
P_U(x) = \frac{1}{\pi}\int_{\mathbb{R}} \frac{y}{(x-\xi)^2+y^2}U(\xi)\,d\xi
\]
is a harmonic function on the upper half plane $\{z\in\mathbb{C}\mid \operatorname{Im}z>0\}$ and it converges to $U(\xi)$ as $z\to\xi$ at a point $\xi$ where $U(\xi)$ is continuous.
\end{exercise}

\begin{solution}
\textbf{Prerequisites:} Harmonic functions, Holomorphic functions, Poisson kernel, and the concept of an Approximate Identity. \\[1em]
\textbf{Intuition:} The Poisson kernel $P_y(x-\xi)$ is a specialized "smearing" function. When we integrate a boundary function $U(\xi)$ against this kernel, it creates a new function $P_U(x,y)$ inside the upper half-plane. This process has two key properties: (1) the result is harmonic (smooth with no peaks/valleys) because the kernel itself is related to a holomorphic function, and (2) as we get closer to the boundary ($y \to 0$), the kernel becomes so "sharp" that it picks up the exact value of $U$ at that point. \\[1em]
\textbf{Steps:}
1. \textbf{Harmonicity:} Recall that a function is harmonic if it is the imaginary (or real) part of a holomorphic function. Let $z = x+iy$. Consider the function $g(z) = \frac{1}{\pi} \frac{i}{z-\xi}$. This is holomorphic in the upper half-plane because the only pole is at $z=\xi$, which is on the real axis. \\[1em]
2. Observe that $\operatorname{Re}(g(z)) = \operatorname{Re} \left( \frac{1}{\pi} \frac{i}{(x-\xi)+iy} \right) = \operatorname{Re} \left( \frac{1}{\pi} \frac{i((x-\xi)-iy)}{(x-\xi)^2+y^2} \right) = \frac{1}{\pi} \frac{y}{(x-\xi)^2+y^2}$.
This is exactly our kernel $P_y(x-\xi)$. Since it is the real part of a holomorphic function, the kernel is harmonic in $(x,y)$. Since $U$ is bounded, we can differentiate under the integral sign, so $P_U$ is also harmonic. \\[1em]
3. \textbf{Approximate Identity Properties:} The kernel $P_y(t) = \frac{1}{\pi} \frac{y}{t^2+y^2}$ satisfies:
(a) $\int_{-\infty}^{\infty} P_y(t) dt = \left[ \frac{1}{\pi} \arctan(t/y) \right]_{-\infty}^{\infty} = \frac{1}{\pi} (\frac{\pi}{2} - (-\frac{\pi}{2})) = 1$.
(b) For any fixed $\delta > 0$, the integral $\int_{|t| \geq \delta} P_y(t) dt \to 0$ as $y \to 0^+$. This means the kernel's mass concentrates at the origin as $y$ gets small. \\[1em]
4. \textbf{Convergence:} Let $\xi_0$ be a point where $U$ is continuous. We want to show $|P_U(x,y) - U(\xi_0)| \to 0$ as $(x,y) \to (\xi_0, 0)$.
Using the fact that the integral of the kernel is $1$, we write:
$|P_U(x,y) - U(\xi_0)| = \left| \int_{\mathbb{R}} P_y(x-\xi) (U(\xi) - U(\xi_0)) d\xi \right|$. \\[1em]
5. Split the integral into two parts: a small interval $|\xi-\xi_0| < \delta$ where $|U(\xi) - U(\xi_0)| < \epsilon$ (by continuity), and the rest of the real line.
The first part is bounded by $\epsilon \int P_y d\xi = \epsilon$.
The second part is bounded by $2 \sup|U| \int_{|\xi-\xi_0| \geq \delta} P_y(x-\xi) d\xi$.
As $y \to 0^+$, the mass of the kernel outside the $\delta$-neighborhood of $x$ vanishes. Since $x \to \xi_0$, eventually the neighborhood $|\xi-x| < \delta/2$ is contained within $|\xi-\xi_0| < \delta$. \\[1em]
6. Thus, the second part also goes to zero, leaving us with a total difference less than any $\epsilon$. This proves that $P_U(x,y) \to U(\xi_0)$.
\end{solution}

\begin{exercise}
Let $f\in L^2(\mathbb{R})$ and let $\hat{f}$ be its Fourier transform. Prove that
\[
\int_{-\infty}^{\infty} x^2|f(x)|^2\,dx \cdot \int_{-\infty}^{\infty} \xi^2|\hat{f}(\xi)|^2\,d\xi \geq \frac{\left(\int_{-\infty}^{\infty}|f(x)|^2\,dx\right)^2}{16\pi^2},
\]
under the condition that the two integrals on the left are bounded.

(Hint: Operators $f(x)\mapsto xf(x)$ and $\hat{f}(\xi)\mapsto\xi\hat{f}(\xi)$ after Fourier transform are non-commuting operators. The inequality is a version of the uncertainty principle.)
\end{exercise}

\begin{solution}
\textbf{Motivation:} This is the Heisenberg Uncertainty Principle. In quantum mechanics, it states that one cannot simultaneously know the position and momentum of a particle. Mathematically, it is a consequence of the non-commutation of the position operator $X$ and the momentum operator $P$. The proof utilizes the Cauchy-Schwarz inequality and integration by parts to relate the variances.

Assume the Fourier transform is defined as $\hat{f}(\xi) = \int_{-\infty}^{\infty} f(x) e^{-2\pi i x \xi} dx$. Then $\widehat{f'}(\xi) = 2\pi i \xi \hat{f}(\xi)$.
By Plancherel's theorem, we have:
\[
\int_{-\infty}^{\infty} \xi^2 |\hat{f}(\xi)|^2 \, d\xi = \frac{1}{4\pi^2} \int_{-\infty}^{\infty} |f'(x)|^2 \, dx.
\]
The inequality we wish to prove is:
\[
\left( \int x^2 |f(x)|^2 \, dx \right) \left( \int |f'(x)|^2 \, dx \right) \geq \frac{1}{4} \left( \int |f(x)|^2 \, dx \right)^2.
\]
Using the Cauchy-Schwarz inequality for the $L^2$ inner product:
\[
\left( \int x^2 |f(x)|^2 \, dx \right) \left( \int |f'(x)|^2 \, dx \right) \geq \left| \int_{-\infty}^{\infty} x f(x) \overline{f'(x)} \, dx \right|^2.
\]
We analyze the integral on the right. Note that:
\[
\int_{-\infty}^{\infty} x f(x) \bar{f}'(x) \, dx = \int_{-\infty}^{\infty} x f(x) d\bar{f}.
\]
Using integration by parts (since the finiteness of the integrals implies $f(x) \to 0$ as $|x| \to \infty$):
\[
\int_{-\infty}^{\infty} x f(x) \bar{f}'(x) \, dx = [x |f(x)|^2]_{-\infty}^{\infty} - \int_{-\infty}^{\infty} (f(x) + x f'(x)) \bar{f}(x) \, dx
\]
\[
= - \int_{-\infty}^{\infty} |f(x)|^2 \, dx - \int_{-\infty}^{\infty} x f'(x) \bar{f}(x) \, dx.
\]
Let $I = \int_{-\infty}^{\infty} x f(x) \bar{f}'(x) \, dx$. Then $I = -\|f\|_2^2 - \bar{I}$.
Thus $I + \bar{I} = -\|f\|_2^2$, which means $2 \operatorname{Re}(I) = -\|f\|_2^2$.
Since $|I|^2 \geq (\operatorname{Re} I)^2$, we have:
\[
|I|^2 \geq \left( \frac{1}{2} \int_{-\infty}^{\infty} |f(x)|^2 \, dx \right)^2 = \frac{1}{4} \left( \int |f|^2 \right)^2.
\]
Substituting back into the original expression:
\[
\int x^2 |f(x)|^2 \, dx \cdot 4\pi^2 \int \xi^2 |\hat{f}(\xi)|^2 \, d\xi \geq \frac{1}{4} \left( \int |f|^2 \right)^2
\]
\[
\int x^2 |f(x)|^2 \, dx \cdot \int \xi^2 |\hat{f}(\xi)|^2 \, d\xi \geq \frac{(\int |f|^2)^2}{16\pi^2}.
\]
The equality holds for Gaussian functions $f(x) = A e^{-ax^2}$.
\end{solution}

\begin{exercise}
Let $\Omega$ be an open domain in the complex plane $\mathbb{C}$. Let $\mathbb{H}$ be the subspace of $L^2(\Omega)$ consisting of holomorphic functions on $\Omega$.

\begin{itemize}
\item[a)] Show that $\mathbb{H}$ is a closed subspace of $L^2(\Omega)$, and hence is a Hilbert space with inner product
\[
(f,g) = \int_{\Omega}f(z)\bar{g}(z)\,dx\,dy, \quad z=x+iy.
\]

\item[b)] If $\{\phi_n\}_{n=0}^{\infty}$ is an orthonormal basis of $\mathbb{H}$, then
\[
\sum_{n=0}^{\infty}|\phi_n(z)|^2 \leq \frac{c^2}{d(z,\Omega^c)^2}, \quad \text{for } z\in\Omega.
\]

\item[c)] The sum
\[
B(z,w) = \sum_{n=0}^{\infty}\phi_n(z)\bar{\phi}_n(w)
\]
converges absolutely for $(z,w)\in\Omega\times\Omega$, and is independent of the choice of the orthonormal basis.
\end{itemize}
\end{exercise}

\begin{solution}
a) By the Mean Value Property for holomorphic functions, for any $z \in \Omega$ and $r < d(z, \partial\Omega)$, $f(z) = \frac{1}{\pi r^2} \int_{D_r(z)} f(w) dA(w)$. By Cauchy-Schwarz, $|f(z)| \leq \frac{1}{\pi r^2} \sqrt{\pi r^2} \|f\|_{L^2(D_r(z))} \leq \frac{1}{\sqrt{\pi} r} \|f\|_{L^2(\Omega)}$.
This implies that $L^2$ convergence in $\mathbb{H}$ implies uniform convergence on compact subsets of $\Omega$. By Weierstrass's theorem, the limit of a sequence of holomorphic functions converging uniformly on compacts is holomorphic. Thus $\mathbb{H}$ is closed in $L^2(\Omega)$.

b) The evaluation functional $L_z(f) = f(z)$ is bounded on $\mathbb{H}$ with $|L_z(f)| \leq \frac{1}{\sqrt{\pi} d(z, \Omega^c)} \|f\|$. By the Riesz Representation Theorem, there exists $K_z \in \mathbb{H}$ such that $f(z) = (f, K_z)$.
Then $\|K_z\|^2 = \sum |\phi_n(z)|^2$. Also $\|K_z\| = \sup_{\|f\|=1} |f(z)| \leq \frac{1}{\sqrt{\pi} d(z, \Omega^c)}$.
Squaring gives $\sum |\phi_n(z)|^2 \leq \frac{1}{\pi d(z, \Omega^c)^2}$. Here $c^2 = 1/\pi$.

c) Absolute convergence: $\sum |\phi_n(z) \bar{\phi}_n(w)| \leq \sqrt{\sum |\phi_n(z)|^2} \sqrt{\sum |\phi_n(w)|^2} < \infty$.
Independence: $B(z,w)$ is the reproducing kernel $K(w,z)$ (in standard notation $f(z) = \int K(z,w) f(w) dA(w)$). It is unique because if there were two such kernels $K_1, K_2$, then $(f, K_{1,z} - K_{2,z}) = 0$ for all $f$, implying $K_{1,z} = K_{2,z}$.
\end{solution}

\subsection{Team}

\begin{exercise}
Calculate the integral:
\[
\int_0^{\infty} \frac{\log x}{1+x^2}\,dx.
\]
\end{exercise}

\begin{solution}
Let $I = \int_0^{\infty} \frac{\log x}{1+x^2} dx$. We use the substitution $x = 1/t$, which gives $dx = -1/t^2 dt$.
As $x \to 0$, $t \to \infty$, and as $x \to \infty$, $t \to 0$.
\[
I = \int_{\infty}^0 \frac{\log(1/t)}{1+(1/t)^2} \left( -\frac{1}{t^2} \right) dt = \int_0^{\infty} \frac{-\log t}{t^2+1} dt = - \int_0^{\infty} \frac{\log t}{1+t^2} dt = -I.
\]
Since $I = -I$, we must have $2I = 0$, which implies $I = 0$.
\end{solution}

\begin{exercise}
Construct an increasing function on $\mathbb{R}$ whose set of discontinuities is precisely $\mathbb{Q}$.
\end{exercise}

\begin{solution}
The set of rational numbers $\mathbb{Q}$ is countable, so we can enumerate it as $\mathbb{Q} = \{q_1, q_2, q_3, \dots\}$.
Define the function $f: \mathbb{R} \to \mathbb{R}$ by:
\[
f(x) = \sum_{q_n < x} \frac{1}{2^n}.
\]
1. $f$ is strictly increasing: If $x < y$, there exists a rational $q_k$ such that $x < q_k < y$. Then $f(y) - f(x) = \sum_{x \leq q_n < y} 2^{-n} \geq 2^{-k} > 0$.
2. Discontinuities: At any $x \in \mathbb{Q}$, say $x = q_k$, the left limit is $f(x^-) = \sum_{q_n < x} 2^{-n}$ and the right limit is $f(x^+) = \sum_{q_n \leq x} 2^{-n}$. The jump is $f(x^+) - f(x^-) = 2^{-k} > 0$. Thus $f$ is discontinuous at every rational.
3. Continuity at irrationals: If $x \notin \mathbb{Q}$, then for any $\epsilon > 0$, there exists $N$ such that $\sum_{n > N} 2^{-n} < \epsilon$. There is a neighborhood of $x$ containing no $q_n$ for $n \leq N$. In this neighborhood, $|f(y) - f(x)| \leq \sum_{n > N} 2^{-n} < \epsilon$. Thus $f$ is continuous at every irrational.
\end{solution}

\begin{exercise}
Prove that any bounded analytic function $F$ over $\{z\mid r<|z|<R\}$ can be written as $F(z)=z^{\alpha}f(z)$, where $f$ is an analytic function over the disk $\{z\mid |z|<R\}$ and $\alpha$ is a constant.
\end{exercise}

\begin{solution}
If $F$ is analytic on the annulus $A = \{r < |z| < R\}$, it has a Laurent series expansion $F(z) = \sum_{n=-\infty}^{\infty} a_n z^n$.
Actually, the statement as written seems to imply something about the behavior at $z=0$ or a specific form. If $r=0$ and $F$ is bounded on $\{0 < |z| < R\}$, then by Riemann's Removable Singularity Theorem, $F$ extends to an analytic function on $|z| < R$. In this case $\alpha = 0$ and $f = F$.
If $r > 0$, the function $F$ is already bounded on the annulus (since the closure is compact). Any $F$ can be written as $F(z) = f(z) + g(1/z)$ where $f$ is analytic on $|z| < R$ and $g$ is analytic on $|z| < 1/r$.
Wait, the common interpretation of this problem in the context of complex analysis contests is usually for non-vanishing functions where $\log F$ is considered, or for $H^\infty$ functions. However, if $F$ is a bounded analytic function on $0 < |z| < R$, then $F$ is analytic on $|z| < R$. 
If the question meant a non-vanishing function $F$, then $F'/F$ is analytic. Let $a_{-1} = \frac{1}{2\pi i} \int_\gamma \frac{F'}{F} dz$ be the winding number $\alpha$. Then $\frac{F'}{F} - \frac{\alpha}{z}$ has a primitive $G(z)$, and $F(z) = e^\alpha e^{G(z)} = z^\alpha f(z)$.
\end{solution}

\begin{exercise}
Let $D\subset\mathbb{R}^n$ be a bounded open set, $f:D\to D$ is a smooth map such that its Jacobian $\left|\frac{\partial f}{\partial x}\right|\equiv 1$, where $D$ denotes the closure of $D$. Prove:

\begin{itemize}
\item[a)] for each small ball $B_{\epsilon}(x)$, there exists a positive integer $k$ such that $f^k(B_{\epsilon}(x))\cap B_{\epsilon}(x)\neq\emptyset$, where $B_{\epsilon}(x)$ denotes the ball centered at $x$ with radius $\epsilon$;
\item[b)] there exists $x\in D$ and a sequence $k_1,k_2,\cdots,k_j,\cdots$ such that $f^{k_j}(x)\to x$ as $k_j\to\infty$.
\end{itemize}
\end{exercise}

\begin{solution}
This is the Poincaré Recurrence Theorem.
a) Since $|\det Jf| = 1$, $f$ is a volume-preserving map. Let $V = \text{vol}(B_{\epsilon}(x))$. If $f^k(B_{\epsilon}) \cap B_{\epsilon} = \emptyset$ for all $k \geq 1$, then $f^k(B_{\epsilon}) \cap f^j(B_{\epsilon}) = \emptyset$ for all $k \neq j$ (by applying $f^{-j}$ and using injectivity/volume preservation). 
Then $\text{vol}(\cup_{k=0}^{\infty} f^k(B_{\epsilon})) = \sum_{k=0}^{\infty} \text{vol}(f^k(B_{\epsilon})) = \sum_{k=0}^{\infty} V = \infty$.
But $\cup f^k(B_{\epsilon}) \subset D$ and $\text{vol}(D) < \infty$, a contradiction. Thus $f^k(B_{\epsilon}) \cap B_{\epsilon} \neq \emptyset$ for some $k$.

b) Let $A$ be the set of points $x \in D$ such that $x$ is recurrent, i.e., $x \in \omega(x)$. The Poincaré Recurrence Theorem states that for any volume-preserving map on a finite volume space, almost every point is recurrent. Thus the set of such $x$ has full measure, and in particular, is non-empty. For such an $x$, there exists a sequence $k_j \to \infty$ such that $f^{k_j}(x) \to x$.
\end{solution}

\begin{exercise}
Let $u$ be a subharmonic function over a domain $\Omega\subset\mathbb{C}$, i.e., it is twice differentiable and $\Delta u = \frac{\partial^2 u}{\partial x^2} + \frac{\partial^2 u}{\partial y^2} \geq 0$. Prove that $u$ achieves its maximum in the interior of $\Omega$ only when $u$ is a constant.
\end{exercise}

\begin{solution}
This is the Strong Maximum Principle for subharmonic functions.
A subharmonic function satisfies the mean value inequality: for any $z_0 \in \Omega$ and $r$ such that $D_r(z_0) \subset \Omega$,
\[
u(z_0) \leq \frac{1}{\pi r^2}\iint_{D_r(z_0)} u(x,y)\,dx\,dy.
\]
Suppose $u$ achieves its maximum $M$ at $z_0 \in \Omega$. Then $u(z) \leq M$ for all $z \in \Omega$.
If $u(z_0) = M$, then $M \leq \frac{1}{\pi r^2}\iint_{D_r(z_0)} u(x,y)\,dx\,dy$.
If there existed $z \in D_r(z_0)$ such that $u(z) < M$, then by continuity, the integral would be strictly less than $M \cdot \pi r^2$, which would imply $M < M$, a contradiction.
Thus $u(z) = M$ for all $z \in D_r(z_0)$.
By a standard connectivity argument, the set $\{z \in \Omega \mid u(z) = M\}$ is both open and closed in $\Omega$. Since $\Omega$ is a domain (connected), $u(z) = M$ for all $z \in \Omega$, i.e., $u$ is constant.
\end{solution}

\begin{exercise}
Suppose that $\phi\in C_0^{\infty}(\mathbb{R}^n)$, $\int_{\mathbb{R}^n}\phi\,dx = 1$. Let $\phi_{\epsilon}(x) = \epsilon^{-n}\phi(x/\epsilon)$, $x\in\mathbb{R}^n$, $\epsilon>0$. Prove that if $f\in L^p(\mathbb{R}^n)$, $1\leq p < \infty$, then $f*\phi_{\epsilon}\to f$ in $L^p(\mathbb{R}^n)$ as $\epsilon\to 0$. It is not true for $p=\infty$.
\end{exercise}

\begin{solution}
First, $\|\phi_\epsilon\|_1 = \int \epsilon^{-n} |\phi(x/\epsilon)| dx = \int |\phi(y)| dy = \|\phi\|_1 = 1$.
Using the property $\int \phi_\epsilon = 1$, we have:
\[
(f * \phi_\epsilon)(x) - f(x) = \int_{\mathbb{R}^n} [f(x-y) - f(x)] \phi_\epsilon(y) dy.
\]
By Minkowski's Inequality for Integrals:
\[
\|f * \phi_\epsilon - f\|_p \leq \int_{\mathbb{R}^n} \|f(\cdot - y) - f(\cdot)\|_p |\phi_\epsilon(y)| dy = \int_{\mathbb{R}^n} \|\tau_y f - f\|_p |\phi_\epsilon(y)| dy.
\]
Let $g(y) = \|\tau_y f - f\|_p$. For $1 \leq p < \infty$, $g(y) \to 0$ as $y \to 0$ (continuity of translation in $L^p$).
Given $\delta > 0$, choose $\eta > 0$ such that $g(y) < \delta$ for $|y| < \eta$.
Split the integral:
\[
\int_{|y|<\eta} g(y) |\phi_\epsilon(y)| dy + \int_{|y|\geq \eta} g(y) |\phi_\epsilon(y)| dy \leq \delta \|\phi\|_1 + 2\|f\|_p \int_{|z| \geq \eta/\epsilon} |\phi(z)| dz.
\]
The second term goes to $0$ as $\epsilon \to 0$ because $\phi \in L^1$. Thus $\limsup \|f * \phi_\epsilon - f\|_p \leq \delta \|\phi\|_1$. Since $\delta$ is arbitrary, the limit is $0$.
For $p=\infty$, let $f = \chi_{(0, \infty)}$. Then $f * \phi_\epsilon$ is a smooth transition, but $\|f * \phi_\epsilon - f\|_\infty \geq 1/2$ (at $x=0$), so it does not converge in $L^\infty$.
\end{solution}
